\chapter{Conclusiones}
\label{cap:capitulo5}

\vspace{1cm}

En este último capítulo se presentan las conclusiones obtenidas durante el desarrollo del proyecto, detallando el grado de cumplimiento de los objetivos y requisitos establecidos inicialmente. Además, se presentan las habilidades y los conocimientos obtenidos a lo largo de la ejecución del trabajo, además de varias líneas futuras que permitirán expandir la funcionalidad del sistema creado.

\section{Objetivos y requisitos cumplidos}

En esta sección se analiza el grado de cumplimiento de los objetivos y requisitos definidos al inicio del proyecto. Para ello, se evalúan los resultados obtenidos durante el desarrollo y se comparan con las especificaciones planteadas.

\subsection{Objetivos}

Se ha conseguido cumplir el objetivo principal, descrito en la \autoref{sec:descripcion}, el cual consistía en el diseño y desarrollo de un sistema inteligente de asistencia sanitaria doméstico, denominado \textit{MyMeds}. El sistema desarrollado integra un dispositivo físico basado en un microcontrolador ESP32 con pantalla táctil y una aplicación móvil Android, permitiendo la gestión de la medicación, la monitorización biomédica básica y la sincronización de información entre ambos componentes.\\

Asimismo, se han alcanzado los distintos subobjetivos planteados al inicio del proyecto:

\begin{enumerate}

\item Se ha diseñado la arquitectura general del sistema, definiendo los distintos módulos hardware y software necesarios para su funcionamiento e integración.\\

\

\item Se ha desarrollado el diseño mecánico del pastillero mediante herramientas de modelado 3D y fabricación mediante impresión 3D, obteniendo un prototipo funcional adaptado a los requisitos del proyecto.

\item Se ha seleccionado e integrado el microcontrolador ESP32 como núcleo principal del sistema, implementando la lógica de control necesaria para la gestión de sensores, almacenamiento, comunicaciones e interfaz de usuario.

\item Se ha diseñado e implementado una interfaz táctil intuitiva que permite gestionar las funcionalidades principales del dispositivo de forma sencilla.

\item Se ha integrado un sensor biomédico MAX30102 que permite la monitorización de la frecuencia cardíaca del usuario, incorporando además mecanismos para el almacenamiento y consulta de las mediciones realizadas.

\item Se ha desarrollado una aplicación móvil Android destinada a la configuración, supervisión y gestión del sistema.

\item Se ha implementado un sistema de comunicación inalámbrica mediante WiFi que permite el intercambio de información entre la aplicación móvil y el dispositivo, incluyendo mecanismos de sincronización bidireccional.

\item Se han realizado diferentes pruebas funcionales destinadas a validar el correcto funcionamiento del sistema desarrollado, verificando aspectos relacionados con la sincronización, la monitorización biomédica, la persistencia de datos y la seguridad.

\end{enumerate}

\subsection{Requisitos}

Además de los objetivos planteados inicialmente, también se han satisfecho los requisitos técnicos establecidos y descritos en la \autoref{sec:requisitos} para el desarrollo del sistema:

\begin{enumerate}

\item Se ha utilizado el microcontrolador ESP32-2432S028 como plataforma principal para el desarrollo del dispositivo.

\item El firmware del dispositivo ha sido implementado en lenguaje C++, mientras que la aplicación móvil ha sido desarrollada utilizando Kotlin y Android Studio.

\item La comunicación entre el dispositivo y la aplicación móvil se realiza mediante conexión WiFi, permitiendo el intercambio bidireccional de información.

\item El sistema desarrollado proporciona una respuesta fluida durante su utilización, permitiendo la actualización de información y la interacción con el usuario en tiempo real.

\item El diseño mecánico ha sido realizado mediante herramientas de modelado 3D y fabricado utilizando impresión 3D con material PLA.

\item Se ha mantenido la filosofía de bajo coste planteada durante la fase de diseño mediante el empleo de componentes electrónicos comerciales ampliamente disponibles.

\item El sistema incorpora mecanismos de almacenamiento persistente para conservar tanto la configuración como el historial de mediciones biomédicas realizadas.

\item El dispositivo dispone de avisos visuales y acústicos destinados a facilitar la correcta administración de la medicación y mejorar la experiencia de uso.

\end{enumerate}

\section{Competencias y conocimientos adquiridos}

Además de las competencias empleadas(\autoref{sec:competencias}) para el desarrollo de este Trabajo Fin de Grado, en su transcurso se han adquirido y reforzado conocimientos en diferentes áreas de la ingeniería, entre los que destacan:

\begin{itemize}

\item Desarrollo de sistemas embebidos basados en microcontroladores ESP32.

\item Integración y utilización de sensores biomédicos para la monitorización de parámetros fisiológicos.

\item Desarrollo de aplicaciones móviles Android mediante Kotlin y Android Studio.

\item Implementación de comunicaciones inalámbricas mediante WiFi para el intercambio de información entre dispositivos.

\item Gestión de almacenamiento persistente de datos en sistemas embebidos y aplicaciones móviles.

\item Diseño mecánico asistido por ordenador mediante herramientas CAD y fabricación mediante impresión 3D.

\item Integración conjunta de componentes hardware y software dentro de un único sistema funcional.

\item Realización de pruebas experimentales y validación de prototipos tecnológicos.

\item Elaboración de documentación técnica y científica mediante \LaTeX.

\end{itemize}

\section{Conclusiones personales}

La realización de este Trabajo Fin de Grado ha supuesto una experiencia especialmente enriquecedora tanto a nivel académico como personal. \textit{MyMeds} ha requerido integrar conocimientos de diferentes áreas de la ingeniería, incluyendo programación embebida, desarrollo de aplicaciones móviles, diseño mecánico e impresión 3D, así como comunicaciones entre dispositivos.\\

Uno de los aspectos más interesantes ha sido la posibilidad de desarrollar un sistema completo, desde el diseño físico del dispensador hasta la implementación de la aplicación móvil y la lógica de funcionamiento del dispositivo. Esta visión global me ha permitido comprender mejor los desafíos asociados al desarrollo de productos reales, donde hardware y software deben trabajar de forma coordinada.\\

La parte más compleja del desarrollo ha sido la sincronización de la información entre la aplicación móvil y el ESP32. Garantizar la consistencia de los datos, evitar conflictos y asegurar el correcto funcionamiento tras reinicios del sistema requirió numerosas pruebas y modificaciones de la arquitectura inicial. Asimismo, la integración de nuevas funcionalidades durante el desarrollo obligó a replantear en varias ocasiones determinadas decisiones de diseño.\\

A pesar de las dificultades encontradas, los resultados obtenidos han sido satisfactorios y han permitido materializar una solución funcional capaz de gestionar la dispensación de medicamentos y el registro de parámetros biomédicos. Además, el proyecto ha contribuido significativamente a ampliar mis conocimientos técnicos y mi capacidad para afrontar proyectos complejos de carácter multidisciplinar.

\section{Líneas futuras}

Aunque los resultados obtenidos demuestran la viabilidad del sistema desarrollado, existen diversas líneas de trabajo que podrían abordarse en futuras versiones con el objetivo de ampliar sus funcionalidades y mejorar sus prestaciones.

\begin{itemize}

\item \textbf{Integración con sistemas sanitarios.} Integrar \textit{MyMeds} con plataformas y sistemas de información sanitaria. Esto permitiría sincronizar automáticamente tratamientos prescritos por profesionales médicos, compartir historiales de medicación y facilitar el seguimiento remoto de pacientes por parte de centros sanitarios o personal autorizado.

\item \textbf{Soporte multiusuario.} Actualmente el sistema está orientado a la gestión de la medicación de un único usuario. Como evolución futura, podría desarrollarse una arquitectura multiusuario que permitiera gestionar varios pacientes desde un mismo dispositivo. Para ello, cada compartimento o conjunto de compartimentos podría asociarse a un usuario diferente, permitiendo almacenar configuraciones, historiales y tratamientos independientes. De esta forma, el sistema podría utilizarse en entornos familiares o residencias de pequeño tamaño.

\item \textbf{Ampliación de la monitorización biomédica.} El sistema podría complementarse con nuevos sensores capaces de registrar parámetros adicionales como la temperatura corporal o la presión arterial, ampliando así las capacidades de seguimiento de la salud del usuario.

\item \textbf{Notificaciones remotas y seguimiento de cuidadores.} También resultaría interesante incorporar mecanismos de notificación dirigidos a familiares o cuidadores. Estas funcionalidades permitirían informar de olvidos en la medicación, incidencias detectadas o valores fisiológicos fuera de los rangos habituales.

\item \textbf{Aplicación de técnicas de inteligencia artificial.} En futuras versiones podrían emplearse técnicas de inteligencia artificial para analizar patrones de comportamiento, detectar posibles incumplimientos terapéuticos y generar recomendaciones personalizadas destinadas a mejorar la adherencia a los tratamientos.

\end{itemize}

En conjunto, estas líneas futuras permitirían evolucionar \textit{MyMeds} desde un prototipo funcional de asistencia sanitaria doméstico hacia un producto más completo, escalable y adaptado a entornos reales de uso.


